\scanpage{172}
\mspage{46}

$\mathrm d_1$) Tout sous-objet d'un objet de $\mathrm{Filt}^{\unsure{\pi|\varpi}}$ est dans $\mathrm{Filt}^{\pi}$

$\mathrm d_2$) Si $Z$ est dans $\mathrm{Filt}^{\pi}$, alors \deleted{tout quotient}
$Z'$ de $Z$ tel que \deleted{l'image} $\mathrm{Im}\,p=Z'$ (où
$p:Z\to Z'$ est le morphisme canonique)
est dans $\mathrm{Filt}^{\pi}$.

% slanted
\margin{Il semblerait qu'on \uline{n'a} pas besoin de $\mathrm d_1$), seulement
de $\mathrm d_2$. Mais en fait, $\mathrm d_1$) doit être conséquence de
$\mathrm d_2$, et $\mathrm d_2$ est conséquence de la \unsure{première partie} de d
(stabilité par \uline{lim} finies) si \unsure{tout} \deleted{\unsure{il}} épimorphisme de $\unsure{\mathcal M|M}$
est strict, donc effectif (si les $\mathrm{Filt}^{\pi}$ sont stables
par \uline{lim} finies \unsure{aussi|-- si} \ldots).}

Bien sûr, \deleted{le plus souvent, si} pratiquement $\mathcal M$ sera
tel que $\mathrm{Filt}^{\pi}$ est stable par
passage aux sous-objets \deleted{quelconques} et aux
quotients d'un objet (pas seulement les
quotients particuliers dont il est question
dans $\mathrm d_2$). D'ailleurs, lorsque \deleted{tout}
bimorphisme (i.e. mono \uline{et} épi) de $\mathcal M$
est iso, alors $\mathrm d_2$ \deleted{implique} \inserted{équivaut} bien à la
stabilité de $\mathrm{Filt}^{\pi}$ par quotients d'objets.

\uline{Proposition} 5.3. Supposons les conditions
a) b) \unsure{d$_0$)} \deleted{\unsure{d)}} vérifiées pour $\mathcal M$. Soit
$\pi$ un cardinal $\geq\pi_0$, et $X$ un objet
de $\mathcal M$. \deleted{Alors} Soit $\mathrm{Ss}_{\pi}(X)$ l'ens.
des sous-objets de $X$ qui sont dans $\mathrm{Filt}^{\pi}$.
Cet ensemble, \unsure{ordonné|muni} par la rel. d'ordre
induite par $\mathrm{Ss}(X)$ est \deleted{filtrant}.
\unsure{gr.|grand} devant $\pi$ (et en particulier filtrant),
\deleted{\unsure{et lorsque} $X$ \unsure{est} \illegible{2} \unsure{isomorphisme}
\unsure{limite inductive} \illegible{1} $X$}
\deleted{\unsure{de plus}} si de plus, les limites
inductives filtrantes de monomorphismes
de $\mathcal M$ sont mono\textsuperscript{*} \inserted{(p. ex. les limites inductives
filtrantes exactes à g.)}, ou si les objets de $\mathrm{Filt}^{\pi}$
sont $\pi$-accessibles, on a
$\underset{\alpha\in\mathrm{Ss}_{\pi}(X)}{\underrightarrow{\mathrm{lim}}} Z_{\alpha}\overset{\sim}{\longrightarrow} X$.
% margin note written sideways
\margin{C'est une \uline{condition \deleted{\unsure{de}} d'exactitude} \unsure{de} $\mathcal M$,
que \unsure{je vais|je veux} noter comme \textcircled{\unsure{e}} \deleted{\unsure{Alors}} les \unsure{mono|monos} de $\mathcal M$ sont stables
par \uline{lim} filtrantes dans $\underline{\mathrm{Fl}}(\mathcal M)$.}
% slanted
\margin{Dire aussi que $\mathrm{Ss}_{\pi}(X)$ est $\pi$-\unsure{complet|adapté|cofiltrant}, et le système
induc\textsuperscript{\unsure{tif}} \unsure{(disp\textsuperscript{f} $\alpha$)} $(Z_\alpha)_\alpha$, \unsure{aussi}.}

% archive sheet number: 171
